\begin{annexitem}{annex:p3.adm.1}{thm:p03-adm-collection}{Collection and replacement}
\begin{proof}
Finite tuples can be coded by iterated Kuratowski pairs. The relation
asserting that a set is the pair of two specified objects is bounded:
its members must be exactly their singleton and unordered pair. The
coordinates are bounded by iterated unions of the code. Pairing, Union
and bounded Collection consequently permit coding or unpacking a finite
tuple of existential witnesses without changing a bounded matrix.

First, Collection for a formula \(\exists z\,\delta(x,y,z)\), with
\(\delta\) bounded, follows directly from bounded Collection applied to
a code of \((y,z)\). A collecting set of such codes gives a bound for
all the first coordinates by taking finitely many unions. This proves
\(\Sigma_1\)-Collection before any appeal to a general normalization
lemma.

Normalize \(\Sigma\)-formulas by induction on their formation.
Conjunctions combine their witness tuples; disjunctions can use a finite
branch tag and the witness for the chosen branch. Existential
quantification adds another coordinate. Bounded existential
quantification is treated the same way, retaining its bound in the
matrix. The only substantial case is bounded universal quantification.
After normalizing its matrix it has the form
\(\forall x\in a\,\exists y\,\delta(x,y)\). Bounded Collection gives the
equivalence
\[
 \forall x\in a\,\exists y\,\delta(x,y)
 \quad\longleftrightarrow\quad
 \exists b\,\forall x\in a\,\exists y\in b\,\delta(x,y).
\]
The right side has bounded matrix. Dummy witnesses handle the empty
bound or unused variables. Thus every \(\Sigma\)-formula normalizes,
and Collection for all of them follows from the first paragraph's
\(\Sigma_1\)-Collection.

Suppose \(\exists y\,\delta(x,y)\) and \(\exists z\,\epsilon(x,z)\)
are complementary on a set \(a\), with bounded matrices. For each
\(x\in a\), collect a tagged witness for one of these alternatives.
There is then a set \(b\) bounding witnesses for the true alternative
at every \(x\). On \(a\), the first alternative is equivalent to the
bounded condition that a tag-zero witness in \(b\) satisfies
\(\delta\): a witness for the wrong alternative cannot satisfy its
matrix because the definitions are complementary. Bounded Separation
therefore gives the required subset of \(a\). This is
\(\Delta_1\)-Separation. The argument is relative to the hypotheses of
complementarity, as asserted.

Finally let \(\sigma(x,y)\) be a \(\Sigma\)-definition of a total
functional relation on \(a\). Collection gives a set \(b\) containing
at least one value for each \(x\). On \(a\times b\), the complement of
\(\sigma(x,y)\) is expressed by
\(\exists z\,(\sigma(x,z)\land z\ne y)\), using totality and
functionality. Both are \(\Sigma\), hence normalizable to
\(\Sigma_1\). Products are sets: for each \(x\in a\) collect the pairs
\(\langle x,y\rangle\), \(y\in b\), and then collect and unite those
rows; bounded Separation removes any extra collected objects.
\(\Delta_1\)-Separation on this product gives the graph, and bounded
Collection and Separation give its range. No choice function for a
nonfunctional relation was used.
\end{proof}
\end{annexitem}

\begin{annexitem}{annex:p3.adm.2}{thm:p03-adm-recursion}{Set recursion}
\begin{proof}
A partial computation is a function \(h\in A\) on an \(R\)-downward
closed subset \(D\in A\) of \(I\), satisfying
\(\gamma(x,h\restriction I_x,h(x))\) for every \(x\in D\).
Downward closure ensures that the displayed restriction has domain
\(I_x\). Function restrictions are formed by bounded Separation.
The assertion that \((D,h)\) is such a computation is \(\Sigma\): the
function and closure conditions are bounded, restrictions can be
existentially recorded, and the remaining universal quantifier ranges
over \(D\). Thus it has a \(\Sigma_1\) form by the preceding lemma.

Any two partial computations agree on their common domain. Otherwise
choose an \(R\)-minimal point of disagreement in the surrounding
universe. Both domains contain all its predecessors. The functions agree
there by minimality, so the uniqueness assumption on \(\gamma\) gives
equal values at the point as well, a contradiction. Here external
well-foundedness is used; it is not replaced by an assertion that \(A\)
decides well-foundedness.

By external well-founded induction on \(R\), each \(x\in I\) belongs
to the domain of a partial computation in \(A\). Assume this for its
predecessors. They form the set \(I_x\in A\). Apply
\(\Sigma_1\)-Collection to their partial-computation witnesses. More
explicitly, normalize the witness condition as \(\exists w\,\delta\)
with bounded \(\delta\), collect tuples \((y,D,h,w)\), and retain by
bounded Separation just tuples with \(y\in I_x\), \(y\in D\), and
\(\delta(y,D,h,w)\). Every predecessor remains covered and every
retained function is a valid partial computation. Collect their domains
and graphs and take their unions. By compatibility these unions give
one partial computation \(h_0\in A\) on a downward closed domain
\(D_0\in A\) containing \(I_x\). If it already contains \(x\), it
suffices. Otherwise use the unique value of
\(\gamma(x,h_0\restriction I_x,z)\) to add \(x\) and its value. The
result is again a partial computation in \(A\).

Now apply the same collection and union argument to all \(x\in I\).
The union is a function \(f\in A\) with domain exactly \(I\), satisfying
the required recursion. Uniqueness follows by the same minimal-point
argument. Finally the statement that the value at \(x\) is \(z\) says
that some partial computation has \(x\) in its domain and value \(z\).
This is a \(\Sigma\)-formula with the displayed fixed data as parameters.
The proof proves correctness on the stipulated well-founded inputs,
not on arbitrary relation codes.
\end{proof}
\end{annexitem}

\begin{annexitem}{annex:p3.adm.3}{thm:p03-adm-syntax}{Syntactic closure}
\begin{proof}
Pairing, Union, Separation and Collection give the elementary set
operations. Recursing on \(\omega\) gives the finite powers of a set;
their union contains all its finite lists. Products exist by
\cref{thm:p03-adm-collection}. A finite computation can be recorded
as a list, so its value relation is \(\Sigma\). Totality and uniqueness
supply a complementary \(\Sigma\) definition. The graph is therefore
\(\Delta_1\).

For \(x\in A\), iterate \(u\mapsto u\cup\bigcup u\) from \(x\)
along \(\omega\). The sequence and its union belong to \(A\) by set
recursion and replacement. That union is transitive, contains \(x\) as
a subset, and is contained in every transitive set containing \(x\)
as a subset, by induction on the iterates. It is \(\mathrm{TC}(x)\).
Existentially recording the iteration gives its \(\Sigma\) graph;
uniqueness and totality give its complementary definition.

The code \(\chi\) and all its subformulas belong to
\(\mathrm{TC}(\{\chi\})\). The immediate-subformula relation
on these codes is a set relation in \(A\) and is well founded: each constituent has smaller
set rank than its containing tagged code. Syntactic predicates and
operations can consequently be evaluated by set recursion over this
closure. At an atomic node, check its finite term and symbol records.
At a connective node, collect the already computed child records;
union their symbol and free-variable sets, and add the connective tag.
At a quantifier node, remove the bound variable from the free-variable
set. A check that a free-variable set is finite is bounded by \(\omega\):
with variables indexed by natural numbers it says that all its indices
are below some \(n\in\omega\). Invalid constructor records are assigned
an invalid flag, so syntactic checking is total on candidate codes.
These steps use bounded operations, finite computations and replacement;
their graphs are \(\Sigma\), and the recursion lemma applies.

For normal form compute simultaneously the positive and negative
translations at each node, using exactly the clauses of
\cref{def:p02-w02c-p02-003-hintikka}. For substitution of a constant,
replace the specified free variable in atomic records; at a quantifier
binding that variable leave its matrix unchanged, and otherwise recurse.
An injective variable-renaming map in \(A\) is applied to the finite
variable records at each node. The fixed shifts used to free a variable
have computable graphs and belong to \(A\), since \(\omega\in A\).
The relational-image construction replaces each finite term record by
its finite defining graph formula, and then recurses through the
connective tree. The function and constant symbols' totality and
uniqueness axioms are obtained from their arity records; replacement
collects these axioms for the set of symbols occurring in an input.
Thus the translations and their auxiliary conjunctions lie in \(A\).
Their semantic equivalences are those already proved in
\cref{thm:p02-w02c-p02-003-nnf,thm:p02-w02c-p02-003-translations}.

All these recursively assigned values are unique. On legal inputs,
existence of a computation giving the specified value is \(\Sigma\),
and existence of a computation giving a different value is its
complement; normalization gives \(\Delta_1\) graphs. The constructor
\(\langle\text{conjunction},\Phi\rangle\), or its disjunction counterpart,
is in \(A\) whenever \(\Phi\in A\). The common finite free-variable
bound is exactly the syntactic requirement for this constructor to be
a formula. No assertion that \(A\) contains every external countable
subset of itself has entered the argument.
\end{proof}
\end{annexitem}

\begin{annexitem}{annex:p3.adm.4}{thm:p03-adm-proof-codes}{Proof codes and their unions}
\begin{proof}
The data of a proof are sets: its node set, child relation, root,
condition labels, rule records and ordinal ranks. Checking that these
are functions or relations on the indicated domains uses bounded
formulas. The predicate that a rank is an ordinal is bounded, and the
strict decrease requirement is bounded by the child relation. Local
rule checking uses syntactic operations from the preceding lemma.
For the disjunction rule, the condition that every \(\chi\in\Phi\)
has a child with label \(\Gamma\cup\{\chi\}\) quantifies only over
\(\Phi\) and the node set. Freshness quantifies over the computed
occurring-symbol sets. Thus checking a proposed code is \(\Delta_1\)
over \(A\): bounded quantification and Boolean combinations preserve
relations having complementary \(\Sigma\) definitions, by normalization.
Existentially quantifying the code gives a \(\Sigma_1\) definition of
refutability.

Soundness follows by external induction on the ordinal ranks. An
initial condition contains contradictory literals. Rules 1--4 add a
sentence already true in a structure satisfying the conclusion. For
rule 5, a structure satisfying the conclusion satisfies one of the
disjuncts, hence one of the required child conditions, contrary to the
induction hypothesis; an empty disjunction has no satisfying structure.
For rule 6, a structure satisfying the conclusion has an existential
witness. Reinterpret the fresh name as that witness. None of the
conclusion's sentences changes truth, whereas the designated child
condition becomes true, again a contradiction. No arbitrary assignment
to a fresh name was treated as a witness without this reinterpretation.

For the many-premise joining assertion, the relation saying that
\(P\) is a proof of \(\Gamma\cup\{\chi\}\) is \(\Delta_1\), and
\(\Phi\in A\). Collection gives a set \(B\in A\) of pairs
\((\chi,P)\) covering every \(\chi\in\Phi\). Remove, by
\(\Delta_1\)-Separation, pairs whose first coordinate is not in \(\Phi\)
or whose second coordinate is not a valid proof with that root label.
The remaining set still covers \(\Phi\). Retain all of its proofs;
there is no need to select one for each \(\chi\).

Take disjoint tagged copies of their node sets and child relations,
using the pair \((\chi,P)\) as tag, and add a new root labelled
\(\Gamma\). Join this root to every copied proof root. Collection and
replacement form the tagged data in \(A\). Their old ordinal ranks form
a set in \(A\); choose for the new root an ordinal strictly above all
of them, obtained by taking their union and then a successor. The old
ranks remain unchanged. This is a proof code in \(A\) applying rule 5.
The empty family gives the permitted zero-premise instance. For any
one-premise rule, the same construction with a single tagged proof and
a larger root rank suffices. This verifies the precise closure used
below.
\end{proof}
\end{annexitem}

\begin{annexitem}{annex:p3.adm.5}{thm:p03-adm-completeness}{Completeness in the fragment}
\begin{proof}
Soundness is the preceding lemma. For the converse let \(\mathcal S\)
consist of the finite conditions with no \(A\)-refutation. It is an
external set; it need not be an element of \(A\). All formulas in it
are countable infinitary formulas, because \(A\) is externally countable.
We check the consistency-property clauses of
\cref{def:p02-w02c-p02-003-consistency}.

Clause (C0) holds because contradictory atomic literals are initial
conditions. Clauses (C1)--(C4) follow by contraposition from the respective
one-premise joining rules: a refutation of the prescribed extension
would give a refutation of the original condition. The extension still
consists of formulas in \(A\) by syntactic closure. For (C5), if every
\(\Gamma\cup\{\chi\}\), \(\chi\in\Phi\), were refutable in \(A\),
the many-premise joining lemma would refute \(\Gamma\). Its hypothesis
\(\Phi\in A\) holds because the tagged disjunction code belongs to the
transitive set \(A\). Therefore at least one disjunct gives a condition
in \(\mathcal S\). For (C6), choose a name from \(C\) not occurring
in \(\Gamma\). This is possible since the occurring-name set is finite.
A refutation of the resulting witness extension would give one of
\(\Gamma\) by rule 6, so this extension also belongs to \(\mathcal S\).

Thus \(\mathcal S\) is a consistency property. The existing model-existence
theorem \cref{thm:p02-w02c-p02-003-model-existence} gives a model of every
one of its conditions. Hence an unsatisfiable condition cannot belong
to \(\mathcal S\), and has an \(A\)-refutation. The external use of
countability and Choice in that theorem is retained; it has not been
misrepresented as an internal choice principle of \(A\).

A sentence \(\eta\in L_A\) is valid exactly when the condition
\(\{\eta^{\ominus}\}\) is unsatisfiable, equivalently when it has an
\(A\)-refutation. Normalization has a \(\Delta_1\) graph and refutability
is \(\Sigma_1\), so this gives the asserted \(\Sigma_1\) validity
predicate. A finite group of premises can be put in the same condition.
For function and constant symbols, validity of \(\eta\) is equivalent
to validity of \(\mathrm{Fn}(\eta)\to\eta^u\), with the relational
translation and totality/uniqueness axioms of the existing interpolation
proof. Both the translation and the conjunction of these axioms belong
to \(A\) by syntactic closure. This reduces that case to the one proved.
\end{proof}
\end{annexitem}

\begin{annexitem}{annex:p3.adm.6}{thm:p03-adm-interpolation}{The admissible interpolant}
\begin{proof}
First use relational signatures. Follow the inseparable-pair construction
in \cref{annex:p02-w02c-p02-003-interpolation}, but require every formula
and every separator to belong to \(A\). Thus a pair
\((\Gamma_1,\Gamma_2)\) consists of finite conditions in the two respective
signatures, and a separator \(\theta\in A\) uses only their common
relation symbols and the finitely many auxiliary names in the pair,
with \(\Gamma_1\models\theta\) and
\(\Gamma_2\models\neg\theta\). The set of inseparable pairs is taken
externally, as in the completeness proof.

Here are the closure checks, including the point at which admissibility
is indispensable. A contradictory pair of literals on the same side
is separated by falsehood or truth. If the conflicting literals are
on opposite sides, their common atomic sentence separates them. Adding
\(c=c\) to one side cannot destroy inseparability: a hypothetical new
separator is already entailed on that side, and, if \(c\) was new,
universal quantification removes it. For atomic substitution using
\(c=d\), if the equality is on the same side it already entails the
substituted atom. If the equality is on the other side, combine a
hypothetical separator with \(\neg(c=d)\) by disjunction. These are
exactly the two cases (C2) of the cited proof and give separators for
the old pair. Adding a conjunct or a universal instance gives a sentence
entailed by its side; a newly used name is again removed by universal
quantification. Adding a fresh existential witness is handled by
existentially quantifying that name out of the hypothetical separator.
The constant-removal argument is the cited proof's (E): injectively
shift the variable stock, replace the name by the newly freed variable,
and quantify. Syntactic closure proves that all these constructions
remain in \(A\), with exactly the indicated common symbols. Reversing
the two sides and negating a separator gives the symmetric checks.

It remains to justify (C5) within \(A\). Suppose
\(\bigvee\Phi\in\Gamma_1\) and each extended pair
\((\Gamma_1\cup\{\chi\},\Gamma_2)\), \(\chi\in\Phi\), has a separator
in \(A\). Existence of a particular separator with its two validity
certificates is a \(\Sigma_1\) condition over \(A\), by
\cref{thm:p03-adm-completeness}. Concretely the certificates refute
\(\Gamma_1\cup\{\chi,\theta^{\ominus}\}\) and
\(\Gamma_2\cup\{\theta^+\}\); symbol restrictions and normal-form
computations are \(\Delta_1\). Since \(\Phi\in A\), Collection gives a
set in \(A\) of such records covering every \(\chi\). Retain only valid
records by \(\Delta_1\)-Separation on their proof certificates, and let
\(\Theta\in A\) be the set of their separator coordinates. Keeping all
valid records, instead of choosing one for each \(\chi\), avoids an
internal choice assumption.

The sentence \(\bigvee\Theta\) belongs to \(A\). All its constituents
are sentences and use only the same common signature and the finite
name set of the old pair: a constituent \(\chi\) of the disjunction
introduces no new name. A model of \(\Gamma_1\) satisfies some
\(\chi\in\Phi\), and therefore at least one separator collected for
that \(\chi\). A model of \(\Gamma_2\) refutes every separator in
\(\Theta\). Thus \(\bigvee\Theta\) separates the old pair. If \(\Phi\)
is empty, the first side is inconsistent and the empty disjunction is
already a separator. Contrapositively an inseparable pair has some
inseparable disjunct extension. This is the missing internal-family
argument; external countable choice alone would not put \(\Theta\) in
\(A\).

The unions of inseparable pairs now satisfy all clauses of the existing
consistency property. If \((\{\varphi^+\},\{\psi^{\ominus}\})\) were
inseparable, model existence would give a model of \(\varphi\land
\neg\psi\), contradicting the hypothesis. It therefore has a separator
\(\theta\in A\). Since the initial pair has no auxiliary names, this
separator uses exactly the permitted common vocabulary and is the
required interpolant.

For signatures with functions and constants use the same relational
images as in Step 2 of
\cref{annex:p02-w02c-p02-003-interpolation}: interpolate between
\(\varphi^u\land\mathrm{Fn}(\varphi)\) and
\(\neg\mathrm{Fn}(\psi)\lor\psi^u\), then translate back. Syntactic
closure keeps both inputs, their interpolant and its inverse translation
in \(A\). The graph axioms' uniqueness establishes the semantic
implications and preserves the intersection of the vocabularies, exactly
as proved there. Thus no broader fragment is needed.
\end{proof}
\end{annexitem}

\begin{annexitem}{annex:p3.adm.7}{thm:p03-adm-beth}{Explicit definitions}
\begin{proof}
Rename \(P\) to a fresh relation \(P'\) in a second copy of
\(\varphi\), and introduce fresh constants \(c_1,\ldots,c_n\). Implicit
uniqueness makes the following implication valid:
\[
 \varphi(P)\land P(\bar c)
 \quad\longrightarrow\quad
 \bigl(\varphi(P')\longrightarrow P'(\bar c)\bigr).
\]
Apply admissible interpolation. Its interpolant \(\eta(\bar c)\in A\)
contains neither \(P\) nor \(P'\), since each appears on only one side.
It contains only the remaining symbols of \(\varphi\) and the new
constants. In a model of \(\varphi(P)\), the first implication gives
\(P(\bar c)\to\eta(\bar c)\). For the reverse implication interpret
\(P'\) as \(P\) and use the second implication. Hence
\(P(\bar c)\leftrightarrow\eta(\bar c)\) in every such model.
Injectively rename variables to free \(n\) variables, and replace the
new constants by them. Both operations preserve membership in \(A\).
The resulting formula explicitly defines \(P\). No existence of a
\(\varphi\)-expansion of every reduct was assumed; the equivalence is
required only in models of \(\varphi\), as the Beth property stipulates.
\end{proof}
\end{annexitem}

\begin{annexitem}{annex:p3.adm.8}{thm:p03-adm-membership-recursion}{Membership recursion in KPU}
\begin{proof}
All arguments in this proof take place in the axiom system. Extend
\cref{thm:p03-adm-collection} to the language with the urelement predicate:
its proofs use only Pairing, Union, bounded Separation and Collection,
so they apply unchanged, with extensionality restricted to sets. An
urelement has no members. For the purposes of taking its members as a
set, use the empty set when the argument is an urelement and the
argument itself otherwise.

First construct transitive closures without using Infinity. A closure
table consists of a transitive set \(t\) and a function \(h\) on \(t\)
satisfying \(\mathrm{Cl}(h,t)\) from the cited definition. In words,
\[
 h(z)=\{v:v\in z\text{ or }(\exists u\in z)\ v\in h(u)\}.
\]
Two closure tables are compatible. Apply foundation to the assertion
that any two tables containing \(z\) agree there. Their domains are
transitive, so the induction hypothesis gives agreement at each member
of \(z\). The displayed equation then gives agreement at \(z\) itself.

Every object \(y\) belongs to the domain of a closure table. Prove this
by the foundation scheme. For each \(u\in y\), assume a table containing
\(u\). The assertion of such a table is \(\Sigma_1\): its local closure
and function checks are bounded. Collection gathers tables covering all
these \(u\); retain valid tuples by bounded Separation and take the
unions of their domains and graphs. Compatibility makes the latter
union a function. The domain is transitive and contains all members of
\(y\). If it already contains \(y\), it suffices. Otherwise add \(y\)
and assign to it the set of its members together with the union of the
sets assigned to those members. These are set operations available by
Collection, replacement and Union. For an urelement the result is empty.
This gives a closure table containing \(y\).

The value assigned to \(y\) is transitive and contains every member of
\(y\). To see transitivity, use foundation induction: if a set \(v\)
is a member of that value, it is either a member of \(y\), in which
case its members lie in the table's value at \(v\), or belongs to the
value at a member of \(y\), which is transitive by induction. In either
case all members of \(v\) belong to the value at \(y\). If a transitive
set \(b\) contains all members of \(y\), the value at \(y\) is contained
in \(b\), again by induction: each member \(u\) of \(y\) is in \(b\),
and its table value is contained in \(b\). Thus the assigned value is
exactly \(\mathrm{TC}(y)\), with the empty value for an urelement.
The table description gives a total functional \(\Sigma_1\) graph for
transitive closure.

Fix \(\bar x,\bar r\). A computation table on a transitive set \(t\)
consists of the functions \(f,h,H\) in \(\Phi_\gamma\), satisfying
\[
 \begin{gathered}
 h(z)=\mathrm{TC}(z),\qquad H(z)=f\restriction\mathrm{TC}(z),\\
 \gamma(\bar x,z,H(z),f(z),\bar r)\qquad(z\in t).
 \end{gathered}
\]
The table predicate is \(\Sigma\), since its outer quantifiers are
bounded and \(\gamma\) is \(\Sigma\).

Any two computation tables agree wherever their domains overlap.
Use foundation induction on the assertion that tables containing \(y\)
agree on \(\{y\}\cup\mathrm{TC}(y)\). For each member \(u\) of \(y\),
both tables contain \(u\) and its transitive closure, and the induction
hypothesis makes them agree there. Since
\[
 \mathrm{TC}(y)=
 \{u:u\in y\}\ \cup\ \bigcup_{u\in y}\mathrm{TC}(u),
\]
their restrictions to \(\mathrm{TC}(y)\) agree. Functionality of
\(\gamma\) makes their values at \(y\) agree as well. This argument
uses the foundation scheme for the displayed quantified assertion, not
an unproved principle selecting a minimal point for the transitive
closure relation.

Every \(y\) belongs to a computation table, again by foundation
induction. Collect computation tables for all \(u\in y\), using
\(\Sigma\)-Collection. Normalize the table predicate and include its
existential certificate in each collected tuple; bounded Separation
then retains only valid tables, still covering every \(u\). Unite their
domains and their \(f,h,H\) graphs. Compatibility gives functions, and
the union domain is transitive. Their restriction equations remain true:
the entire \(\mathrm{TC}(z)\) was already in the domain of every table
containing \(z\), so a compatible union adds no new argument to that
restriction. If the union already contains \(y\), the table suffices.
Otherwise it contains \(\mathrm{TC}(y)\). Form the restriction of its
\(f\) to that set, use the assumed total operation \(G\) to obtain the
value at \(y\), and add \(y\) and the corresponding entries in
\(f,h,H\). The result is a computation table containing \(y\).

It follows that \(\Phi_\gamma(\bar x,y,w,\bar r)\) has exactly one
value \(w\) for every \(\bar x,y\). Define \(F\) by this formula.
Compatibility identifies the restriction recorded by any computation
table with \(F_{\bar x}\restriction\mathrm{TC}(y)\); the latter is a set
by functional \(\Sigma\)-replacement. Hence the local equation in the
table is exactly the required recursion equation. Any other definable
operation satisfying that equation agrees by the same foundation
induction, with the same totality and set-restriction hypotheses.
Finally \(\Phi_\gamma\) is \(\Sigma\) by the bounded formulas displayed
in the earlier definition and closure of \(\Sigma\) under bounded and
existential quantification. None of these internal derivations used an
axiom of infinity or the external well-foundedness of a model.
\end{proof}
\end{annexitem}

\begin{annexitem}{annex:p3.adm.9}{thm:p03-adm-truncation}{The well-founded part}
\begin{proof}
Write \(E\) for the model's membership relation. Let \(W\) consist
of the urelements and all objects whose downward \(E\)-closure is
externally well founded. It is an initial substructure: every member
of an object in \(W\) is in \(W\). Conversely an object all of whose
members are in \(W\) is itself in \(W\). To verify the converse, any
nonempty subset of its downward closure either consists only of the
object or meets the downward closure of one of its members; external
well-foundedness there supplies a minimal member. Equivalently under
the ambient Choice assumption, an infinite descending sequence from
the object would continue from one fixed member, contradicting that
member's well-foundedness.

Membership recursion in the model gives its internal ordinal rank
function: an urelement has rank zero, and a set has rank
\[
 \rho(x)=\sup\{\rho(y)+1:yEx\}.
\]
To apply \cref{thm:p03-adm-membership-recursion}, the operation takes
this union from the recorded predecessor values, with a fixed default
value on arguments that do not code the required function. Collection
forms the set of successor ranks. Foundation induction proves that
all outputs are ordinals and that the operation agrees with the identity
on ordinals. Its graph is total functional \(\Sigma\), hence
\(\Delta_1\) by the collection lemma.

An object is in \(W\) exactly when its internal rank is in \(W\).
For the forward direction, external induction on \(E\) shows that
every member of \(\rho(x)\) belongs to one of the well-founded successor
ranks \(\rho(y)+1\), and therefore is well founded. Thus \(\rho(x)\)
is well founded. For the converse, a decreasing \(E\)-chain gives a
strictly decreasing chain of internal ranks; below a well-founded ordinal
this is impossible. In minimum-element language, choose a point with
least rank in any nonempty subset of the downward closure. This proves
the equivalence without assuming that all ordinals of the model are
externally well founded.

Bounded formulas with parameters in \(W\) are absolute between \(W\)
and \(\mathfrak B\), by
\cref{thm:p03-w03-p03-039-nadel-persistence}, or directly by induction
on their bounded quantifiers. Empty set, Pairing and Union restrict to
\(W\); every object they construct has well-founded members.
Extensionality also restricts, since no members of a retained set are
lost. A boundedly separated subset of a retained set has only retained
members and so is retained. External well-foundedness supplies every
instance of the foundation scheme in \(W\). It remains to verify
Collection; it does not follow just by restricting an arbitrary
collecting set from the larger model.

Let \(a,\bar p\in W\), let \(\delta\) be bounded, and suppose
\[
 W\models\forall x\in a\,\exists y\,\delta(x,y,\bar p).
\]
For each \(xEa\) there is a witness \(y\in W\). Absoluteness makes it
a witness in \(\mathfrak B\), with well-founded rank. Inside
\(\mathfrak B\), Collection and ranks supply an ordinal bound for a
set of witnesses. Thus some ordinal satisfies
\[
 P(\alpha):\quad
 \forall x\in a\,\exists y\,
       \bigl(\rho(y)<\alpha\land\delta(x,y,\bar p)\bigr).
\]
Here the rank notation abbreviates its \(\Delta_1\) graph. The internal
foundation scheme gives a least ordinal \(\alpha\) satisfying
\(P\). This uses foundation for the definable class of ordinals
satisfying \(P\), not Separation for an arbitrary unbounded predicate.

This least bound belongs to \(W\). Otherwise it has an ordinal member
\(\beta\notin W\), since an ordinal whose members are all well founded
is well founded. Every well-founded ordinal of the model is strictly
below \(\beta\): internal ordinal comparability rules out equality,
and an ordinal below a well-founded ordinal is well founded. The
already chosen external witnesses in \(W\) therefore show
\(\mathfrak B\models P(\beta)\). The bounded quantifier over \(a\)
ranges over exactly the same objects in both structures. This
contradicts the minimality of \(\alpha\).

Apply internal \(\Sigma\)-Collection once more to obtain a set \(b\)
containing bounded-rank witnesses for every \(xEa\). Use internal
\(\Delta_1\)-Separation to retain exactly its elements satisfying
\[
 \rho(y)<\alpha\quad\land\quad
          (\exists x\in a)\,\delta(x,y,\bar p).
\]
This is a \(\Delta_1\) condition because rank is total and
\(\Delta_1\), and the other quantifier is bounded. The resulting set
\(b'\) still contains a witness for each \(xEa\). Each element has a
well-founded rank, so is in \(W\); hence \(b'\in W\). Absoluteness
now makes \(b'\) the required collecting set in \(W\).

Finally collapse \(W\) externally, fixing each urelement and defining
\(c(x)=\{c(y):yEx\}\) for each set. Well-founded recursion defines this
map. Induction and extensionality make it injective on sets; urelements
are kept in their disjoint designated sort. Its image is transitive over
the urelements and it is an isomorphism for membership and the inherited
relations. The preceding verifications show that this image satisfies
\(\mathrm{KPU}\), hence is admissible. A set whose members are exactly
all the urelements has only well-founded members, so it lies in \(W\).
If the model's internal \(\omega\) is standard, all its members are
exactly the standard finite ordinals, and its collapse is the actual
\(\omega\). The qualification excludes the nonstandard case.
\end{proof}
\end{annexitem}
